\begin{tablalafrox}{Diccionario del Producto 6: Solución certificada y documentada. Fuente: según cita por cuenta en cada encabezado.}{tab:T14-dic-prod-6}
  {L{1.2cm} Y Y Y C{0.9cm} L{2.3cm}}
  {\cab{Código} & \cab{Paquete de trabajo} & \cab{Entregable} & \cab{Criterio de aceptación} & \cab{Resp.} & \cab{Período}}
  \multicolumn{6}{>{\columncolor{lafroxBg}\RaggedRight\arraybackslash}p{\dimexpr\textwidth-2\tabcolsep\relax}}{\textbf{Cuenta 6.1 — Pruebas de integración y regresión.} Esta cuenta reúne las pruebas de integración, regresión e idempotencia en QA de la Etapa 1 para el H4. En la Etapa 2, las pruebas para el H9 verifican la integración y que la Etapa 1 en producción no se degrade. {\color{lafroxGris}Fuente: (Bases Administrativas, arts. 17.2–17.6 y Formulario T-14, p. 63); (Bases Técnicas Transversales, secciones 17.2–17.6, pp. 32–34).}} \tabularnewline*
  6.1.1 & Pruebas de integración, regresión e idempotencia en QA (H4) & Informe de pruebas. & Todos los flujos se ejecutan sin error, y no hay regresiones. & CAL & E1 · Ejecución: mes 9. Hito H4: mes 10. \\
  6.1.2 & Pruebas de integración y regresión sin afectar la Etapa 1 en producción (H9) & Informe de pruebas. & No hay regresiones en la Etapa 1. & CAL & E2 · Ejecución: mes 16. Hito H9 del Formulario E-25: mes 17. \\
  \midrule
  \multicolumn{6}{>{\columncolor{lafroxBg}\RaggedRight\arraybackslash}p{\dimexpr\textwidth-2\tabcolsep\relax}}{\textbf{Cuenta 6.2 — Pruebas de aceptación, accesibilidad y condiciones físicas.} Esta cuenta reúne las pruebas de aceptación y accesibilidad. En la Etapa 1 participan preventistas, conductores y personal de bodega; se verifican cinco condiciones físicas y una desconexión de 14 horas. En la Etapa 2 participan cadenas, transportistas, Comercial y Finanzas, con accesibilidad WCAG 2.2 AA en los portales. {\color{lafroxGris}Fuente: (Bases Administrativas, arts. 17.2–17.6 y Formulario T-14, p. 63); (Bases Técnicas Transversales, secciones 17.2–17.6, pp. 32–34).}} \tabularnewline*
  6.2.1 & Pruebas de aceptación y de accesibilidad con preventistas, conductores y bodega & Casos de aceptación firmados. & Los casos están firmados por la Contraparte Técnica, y la accesibilidad cumple WCAG 2.2 AA. & CAL & E1 · Antes del H5. \\
  6.2.2 & Cinco condiciones físicas: −22 °C con guantes, lluvia, calor superior a 34 °C, sol directo y uso con una mano; desconexión separada de 14 horas & Informe de las cinco condiciones y de la desconexión. & Se verifican las cinco condiciones, sin pérdida ni duplicación de registros y con sincronización en 10 minutos o menos. & CAL & E1 · Antes del H5. \\
  6.2.3 & Pruebas de aceptación y de accesibilidad con cadenas, transportistas, Comercial y Finanzas & Casos firmados. & Los casos están firmados por la Contraparte Técnica, y los portales cumplen WCAG 2.2 AA. & CAL & E2 · Antes del H10. \\
  \midrule
  \multicolumn{6}{>{\columncolor{lafroxBg}\RaggedRight\arraybackslash}p{\dimexpr\textwidth-2\tabcolsep\relax}}{\textbf{Cuenta 6.3 — Pruebas de carga y resiliencia.} Esta cuenta reúne las pruebas de carga y resiliencia. En la Etapa 1 se verifican los tiempos a 1,5 veces el peak de septiembre y la recuperación sin intervención; en la Etapa 2 se prueba la carga con ambas etapas activas. {\color{lafroxGris}Fuente: (Bases Administrativas, arts. 17.2–17.6 y Formulario T-14, p. 63); (Bases Técnicas Transversales, secciones 17.2–17.6, pp. 32–34).}} \tabularnewline*
  6.3.1 & Pruebas de carga con el peak de septiembre y de resiliencia & Informe de carga y resiliencia. & Se cumplen los tiempos exigidos a 1,5 veces el peak, y la solución se recupera sin intervención. & CAL & E1 · Antes del H5. \\
  6.3.2 & Pruebas de carga con ambas etapas activas y de resiliencia & Informe. & Se cumplen los tiempos a 1,5 veces el peak. & CAL & E2 · Antes del H10. \\
  \midrule
  \multicolumn{6}{>{\columncolor{lafroxBg}\RaggedRight\arraybackslash}p{\dimexpr\textwidth-2\tabcolsep\relax}}{\textbf{Cuenta 6.4 — Pruebas de recuperación ante desastres.} Esta cuenta reúne las pruebas de recuperación ante desastres mediante conmutación real. La Etapa 1 mide RTO de hasta cuatro horas y RPO de hasta 15 minutos y verifica enlaces, alarmas y conciliación de retorno; la Etapa 2 repite la prueba con ambas etapas activas. {\color{lafroxGris}Fuente: (Bases Administrativas, arts. 17.2–17.6 y Formulario T-14, p. 63); (Bases Técnicas Transversales, secciones 17.2–17.6, pp. 32–34).}} \tabularnewline*
  6.4.1 & Prueba de recuperación ante desastres con conmutación real & Informe AL-DR-01, mediciones y registro de correcciones y repetición. & RTO \ensuremath{\leq}4 h y RPO \ensuremath{\leq}15 min medidos en conmutación real; tres enlaces, alarmas 5/15 min, retorno y conciliación verificados. Las brechas impiden la certificación hasta su cierre. & CAL & E1 · Antes del H5; preparación, ejecución, correcciones y cierre dentro de la ventana de T-15. \\
  6.4.2 & Prueba de recuperación ante desastres con ambas etapas & Informe AL-DR-01 y correcciones con ambas etapas activas. & Conmutación real con RTO \ensuremath{\leq}4 h, RPO \ensuremath{\leq}15 min, alarmas 5/15 min y conciliación de retorno, sin degradar E1. Se repite lo afectado tras corregir brechas. & CAL & E2 · Antes del H10, dentro de la ventana de T-15; preparación, ensayo y cierre incluidos. \\
  \midrule
  \multicolumn{6}{>{\columncolor{lafroxBg}\RaggedRight\arraybackslash}p{\dimexpr\textwidth-2\tabcolsep\relax}}{\textbf{Cuenta 6.5 — Pruebas de seguridad ofensiva.} Esta cuenta reúne las pruebas de seguridad ofensiva realizadas por un tercero independiente para cada etapa. Se entregan los informes y la evidencia de corrección, sin hallazgos críticos ni altos abiertos. {\color{lafroxGris}Fuente: (Bases Administrativas, arts. 17.2–17.6 y Formulario T-14, p. 63); (Bases Técnicas Transversales, secciones 17.2–17.6, pp. 32–34).}} \tabularnewline*
  6.5.1 & Prueba de seguridad ofensiva de la Etapa 1, por un tercero independiente & Informe íntegro del tercero y evidencia de las correcciones. & No quedan hallazgos críticos ni altos abiertos. & SEG & E1 · Antes del H5. \\
  6.5.2 & Prueba de seguridad ofensiva de la Etapa 2, por un tercero independiente & Informe y correcciones. & No quedan hallazgos críticos ni altos abiertos. & SEG & E2 · Antes del H10. \\
  \midrule
  \multicolumn{6}{>{\columncolor{lafroxBg}\RaggedRight\arraybackslash}p{\dimexpr\textwidth-2\tabcolsep\relax}}{\textbf{Cuenta 6.6 — Despliegue sin interrupción demostrado.} Esta cuenta reúne las demostraciones de despliegue y reversión en Preproducción. La primera verifica la Etapa 1 sin interrupción ni intervención manual; la segunda demuestra el procedimiento sin interrumpir ninguna de las dos etapas. {\color{lafroxGris}Fuente: (Bases Administrativas, arts. 17.2–17.6 y Formulario T-14, p. 63); (Bases Técnicas Transversales, secciones 17.2–17.6, pp. 32–34).}} \tabularnewline*
  6.6.1 & Demostración en Preproducción del despliegue sin interrupción de la Etapa 1 & Informe de la demostración. & La versión se despliega y se revierte en Preproducción sin interrupción y sin intervención manual. & SRE & E1 · Antes del H5 y antes de cada paso a producción. \\
  6.6.2 & Demostración en Preproducción del despliegue sin interrupción de la Etapa 2 & Informe de la demostración. & La versión se despliega y se revierte sin interrumpir ninguna de las dos etapas. & SRE & E2 · Antes del H10. \\
  \midrule
  \multicolumn{6}{>{\columncolor{lafroxBg}\RaggedRight\arraybackslash}p{\dimexpr\textwidth-2\tabcolsep\relax}}{\textbf{Cuenta 6.7 — Actas de certificación.} Esta cuenta reúne las actas y la evidencia de certificación de cada etapa: H5 para la Etapa 1 y H10 para la Etapa 2, cuyo cierre incluye el desarrollo. {\color{lafroxGris}Fuente: (Bases Administrativas, arts. 17.2–17.6 y Formulario T-14, p. 63); (Bases Técnicas Transversales, secciones 17.2–17.6, pp. 32–34).}} \tabularnewline*
  6.7.1 & Acta de certificación de la Etapa 1 (H5) & Acta del H5 con su expediente de evidencia. & Las pruebas 6.2.1, 6.2.2, 6.3.1, 6.4.1 y 6.5.1 y la demostración del despliegue sin interrupción 6.6.1 están aprobadas, y el acta está firmada. & CAL & E1 · Ejecución: meses 10 y 11; el acta es documental y no interviene sistemas. Hito H5 del Formulario E-25: mes 12. \\
  6.7.2 & Acta de certificación de la Etapa 2 y cierre del desarrollo (H10) & Acta del H10. & Las pruebas 6.2.3, 6.3.2, 6.4.2 y 6.5.2 y la demostración del despliegue sin interrupción 6.6.2 están aprobadas, y el acta está firmada. & CAL & E2 · Ejecución: meses 16 y 17. Hito H10 del Formulario E-25: mes 18. \\
  \midrule
  \multicolumn{6}{>{\columncolor{lafroxBg}\RaggedRight\arraybackslash}p{\dimexpr\textwidth-2\tabcolsep\relax}}{\textbf{Cuenta 6.8 — Documentación técnica, código y custodia.} Esta cuenta reúne la documentación, el código fuente, la custodia y los elementos necesarios para entender, mantener y traspasar la solución (BA, art. 77, p. 39). Cada entrega liberada se acompañará al CLIENTE con su SBOM de 6.8.4 en formato CycloneDX o SPDX, identificando los componentes, sus identificadores de edición y licencias (BTT, cap. 11, RT-11.23, p. 24). {\color{lafroxGris}Fuente: (Bases Administrativas, arts. 17.2–17.6 y Formulario T-14, p. 63); (Bases Técnicas Transversales, secciones 17.2–17.6, pp. 32–34).}} \tabularnewline*
  6.8.1 & Documentación técnica, funcional, de arquitectura, de operación y de seguridad & Repositorio documental con versiones. & Cada versión en producción tiene su documentación en la misma versión. & ARQ & Ambas · Se revisa antes del H7 y del H12. \\
  6.8.2 & Base de conocimiento de incidentes y decisiones de diseño & Base de conocimiento. & Cada incidente cerrado y cada decisión de arquitectura tienen su entrada. Desde la marcha blanca de la Etapa 1 se miden cada mes su uso y la utilidad declarada por los usuarios; 10.4.3 continúa la medición en la Operación. & SRE & Ambas · Desde el H4. \\
  6.8.3 & Estándares de codificación y documentación de interfaces en OpenAPI y AsyncAPI & Documento de estándares y documentación de interfaces publicada. & Cada interfaz tiene su documentación generada y versionada, y la cadena de integración verifica los estándares. & DES & Ambas · Desde el mes 4, y actualizada en cada versión. \\
  6.8.4 & Inventario de componentes (SBOM) con sus licencias de código abierto & Inventario de componentes actualizado en cada versión. & Cada componente tiene su versión y su licencia, y ninguna licencia es incompatible con el uso del CLIENTE. & DES & Ambas · Desde el mes 4, y en cada versión. \\
  6.8.5 & Depósito del código fuente en el repositorio del CLIENTE en cada liberación & Código depositado en cada liberación. & Cada versión en producción está en el repositorio del CLIENTE, con la misma etiqueta de versión. & DES & Ambas · En cada liberación, desde la primera. \\
  6.8.6 & Entrega del código fuente, la infraestructura como código y los procedimientos de despliegue & Repositorios entregados. & El equipo de TI de Puelche reconstruye y despliega una versión en un ambiente de prueba siguiendo solo lo entregado. & DES & E2 · Mes 21 (H12). \\
  6.8.7 & Contrato de custodia de fuentes con un tercero independiente & Contrato de custodia firmado y primer depósito. & El contrato incluye las cláusulas de liberación ante insolvencia o incumplimiento grave, y el primer depósito está hecho. & JP & Ambas · Antes del H4. \\
\end{tablalafrox}
